\exercisesection{5}

¶ 1. Let A be an integral domain, K its field of fractions and T a linear topology on A (Chapter III, § 2, Exercise 21).

\begin{enumerate}
\def\labelenumi{(\alph{enumi})}
\item
  For the neighbourhoods of 0 under \(\mathcal{T}\) to constitute a fundamental system of neighbourhoods of 0 for some topology \(\mathcal{T}_{\mathrm{K}}\) compatible with the ring structure on K, it is necessary and sufficient that \(\mathcal{T}\) be the topology \(\mathcal{T}_{u}(\mathrm{A})\) (Chapter III, §2, Exercise 24); then A is a bounded subset with respect to \(\mathcal{T}_{\mathrm{K}}\) and \(\mathcal{T}_{\mathrm{K}}\) is a locally bounded Hausdorff topology (General Topology, Chapter III, §6, Exercises 12 and 20 (e)). For K to be complete (resp. linearly compact. resp. strictly linearly compact (Chapter III, §2, Exercise 21)) with \(\mathcal{T}_{\mathrm{K}}\) , it is necessary and sufficient that A be so with \(\mathcal{T}\) .
\item
  For the topology \(\mathcal{T}_{\mathbf{K}}\) (where \(\mathcal{T} = \mathcal{T}_u(\mathbf{A})\) ) to be compatible with the field structure on \(\mathbf{K}\) , it is necessary and sufficient that the Jacobson radical \(\Re(\mathbf{A})\) of \(\mathbf{A}\) be \(\neq 0\) .
\end{enumerate}

¶ 2. Let K be a (commutative) field and T a non-discrete Hausdorff topology on K compatible with the ring structure on K. For the topology T to be defined by a valuation on K or an absolute value on K, it is necessary and sufficient that T be locally retrobounded (General Topology, Chapter III, § 6, Exercise 22. If there exist in K topologically nilpotent elements, use Exercise 22 (d) of General Topology, Chapter III, § 6 and Exercise 13 of General Topology, Chapter IX, § 3. In the opposite case, use Exercise 22 (f) of General Topology, Chapter III, § 6).

¶ 3. Let A be a Noetherian integral domain, K its field of fractions, \(\mathcal{T}_{u}\) the topology \(\mathcal{T}_{u}(\mathrm{A})\) and \(\mathcal{T}_{\mathrm{K}}\) the corresponding topology on K (Exercise 1).

\begin{enumerate}
\def\labelenumi{(\alph{enumi})}
\item
  If A is a Zariski ring with Jacobson radical \(\mathfrak{r} \neq 0\) and A is complete with the r-adic topology, show that A is complete with the topology \(\mathcal{T}_u\) (General Topology, Chapter III, § 3, no. 5, Corollary 2 to Proposition 9).
\item
  Suppose that \(\mathcal{T}_{\mathbf{K}}\) is not discrete and is defined by a valuation \(v\) on \(\mathbf{K}\) . Show that \(v\) is a discrete valuation and that \(\mathbf{A}\) is an open subring of the ring \(\mathbf{V}\) of \(v\) ; obtain the converse. (Using the hypothesis, which implies \(\mathbf{A} \neq \mathbf{K}\) , and Proposition 1 of §4, no. 1, we may assume that \(\mathbf{A} \subset \mathbf{V}\) . Using the fact that \(\mathbf{A}\) is open under \(\mathcal{T}_{\mathbf{K}}\) , show that \(\mathbf{V}\) is a finitely generated A-module and conclude using §3, no. 5, Proposition 9. For the converse, observe that \(\mathbf{A}\) is a local ring of which the maximal ideal \(m\) and (0) are the only prime ideals and in which every ideal \(a \neq 0\) is therefore \(m\) -primary). The integral closure of \(\mathbf{A}\) is then \(\mathbf{V}\) . Give an example where \(\mathbf{A} \neq \mathbf{V}\) (take \(\mathbf{V}\) to be a ring of formal power series \(k[[T]]\) , where \(k\) is a field).
\item
  Deduce an example of a non-discrete complete topological field whose topology is not locally retrobounded (use (b) and Exercise 2).
\end{enumerate}

\begin{enumerate}
\def\labelenumi{\arabic{enumi}.}
\setcounter{enumi}{3}
\tightlist
\item
  Let \(v\) be a valuation on a field \(K\) , A the ring of \(v\) and \(m\) its maximal ideal.
\end{enumerate}

\begin{enumerate}
\def\labelenumi{(\alph{enumi})}
\item
  For the topology defined by v on A to be identical with the m-adic topology, it is necessary and sufficient that A be a field or a discrete valuation ring.
\item
  For the topology defined by \(v\) to make \(A\) into a strictly linearly compact ring (Chapter III, §2, Exercise 21), it is necessary and sufficient that \(A\) be a field or a discrete valuation ring (use (a) and Chapter III, §2, Exercise 22 (a)).
\end{enumerate}

\begin{enumerate}
\def\labelenumi{\arabic{enumi}.}
\setcounter{enumi}{4}
\tightlist
\item
  \begin{enumerate}
  \def\labelenumii{(\alph{enumii})}
  \tightlist
  \item
    Let K be a field, v a valuation on K and \(\Gamma\) the order group of v. For K to be linearly compact (Chapter III, §2, Exercise 15) with the topology \(T_{v}\) , it is necessary and sufficient that, for every well-ordered subset B of \(\Gamma\) and every family \((a_{\beta})_{\beta\in\mathcal{B}}\) of elements of K such that, for \(\lambda<\mu<v\) ,
  \end{enumerate}
\end{enumerate}

\[
v (a _ {\lambda} - a _ {\mu}) <   v (a _ {\mu} - a _ {\nu}),
\]

there exists an element \(a \in \mathbf{K}\) such that \(v(a - a_{\lambda}) = v(a_{\lambda} - a_{\mu})\) for every ordered pair of indices \(\lambda, \mu\) such that \(\mu > \lambda\) . (Use Exercise 4 of Set Theory,

Chapter III, § 2). If this is so, the ring A of the valuation \(v\) is also linearly compact with the discrete topology.

\begin{enumerate}
\def\labelenumi{(\alph{enumi})}
\setcounter{enumi}{1}
\tightlist
\item
  Show that the field \(S(\Gamma, k) = K\) defined in §3, Exercise 2 is linearly compact with \(T_{v}\) . Take \(\Gamma\) to be the totally ordered group Q of rational numbers and consider the subring \(K_{0}\) of K consisting of the \(x = (x_{\alpha})\) such that the set of \(\alpha \in Q\) for which \(x_{\alpha} \neq 0\) is finite or is the set of points of a strictly increasing sequence tending to \(+\infty\) . Show that \(K_{0}\) is a field which is complete with respect to the valuation \(v_{0}\) induced by the canonical valuation v on K and that \(v_{0}\) has the same value group and residue field as v but that \(K_{0}\) is not linearly compact (cf.~§10, Exercise 2).
\end{enumerate}

¶ 6. (a) Let A be a valuation ring, M an A-module and M' a submodule of M. Show that for M' to be a pure submodule of M (Chapter I, § 2, Exercise 24), it is necessary and sufficient that, for all \(\alpha \in A\) , M' ∩ ( \(\alpha\) M) = \(\alpha\) M'.

\begin{enumerate}
\def\labelenumi{(\alph{enumi})}
\setcounter{enumi}{1}
\item
  Let \(\mathbf{M}'\) be a pure submodule of \(\mathbf{M}\) such that, if \(\mathbf{M}'\) is given the topology for which the \(\alpha\mathbf{M}'\) (where \(\alpha \in \mathbf{A}, \alpha \neq 0\) ) form a fundamental system of neighbourhoods of \(0\) , \(\mathbf{M}'\) is linearly compact. Show that for every element \(x \in \mathbf{M}\) there exists \(y' \in \mathbf{M}'\) such that \(x' = x + y'\) has the following property: for all \(\alpha \in \mathbf{A}\) such that \(x + \mathbf{M}' \in \alpha(\mathbf{M} / \mathbf{M}')\) , \(x' \in \alpha\mathbf{M}\) . (For each of the \(\alpha \in \mathbf{A}\) such that \(x + \mathbf{M}' \in \alpha(\mathbf{M} / \mathbf{M}')\) , consider the subset \(\mathrm{S}_{\alpha}\) of \(\mathbf{M}'\) consisting of the \(y_{\alpha}'\) such that \(x + y_{\alpha}' \in \alpha\mathbf{M}\) .)
\item
  Suppose that A is a linearly compact valuation ring; let K be the field of fractions of A. Show that, if M is a torsion-free A-module such that \(\mathbf{M}_{(\mathbf{K})}\) admits a countable basis, M is the direct sum of a countable family of A-modules of rank 1. (Consider M as the union of an increasing sequence \((\mathbf{M}_{i}^{\prime})\) of pure submodules such that \(M_{i}^{\prime}\) is of rank i and for each i apply (b) to \(M_{i}^{\prime}\) and its submodule \(M_{i-1}^{\prime}\) .
\item
  With the same hypotheses on A and M as in (c), let N be a pure submodule of M of finite rank. Show that N is a direct factor of M (observe that every A-module which is the direct sum of a finite number of A-modules of rank 1 is linearly compact and use (b)).
\item
  With the same hypotheses on M and \(M'\) as in (b), show that for all \(x \in M\) there exists \(y' \in M'\) such that the annihilator of \(x' = x + y'\) is equal to the annihilator of the element \(x + M' \in M/M'\) . (Let a be the annihilator of \(x + M'\) ; for all \(\alpha \in a\) , let \(T_{\alpha}\) be the subset of \(M'\) consisting of the \(y'\) such that \(\alpha(x + y') = 0\) ; show that the intersection of the \(T_{\alpha}\) is non-empty.)
\item
  Suppose that A is a valuation ring such that for every ideal \(a \neq 0\) the A-module A/a is linearly compact (with the discrete topology). Show that every finitely generated torsion A-module M is the direct sum of a finite number of monogenous A-modules. (If \((z_{i})_{1 \leqslant i \leqslant n}\) is a system of generators of M, argue by induction on n, considering one of the \(z_{i}\) whose annihilator is the least and noting that the submodule of M which it generates is pure.)
\end{enumerate}

\begin{enumerate}
\def\labelenumi{\arabic{enumi}.}
\setcounter{enumi}{6}
\tightlist
\item
  Let K be a field and v a discrete valuation on K such that K is complete with \(T_{v}\) : let A and p denote the ring and ideal of v and U = A - p the set of invertible elements of A. Let u be an isomorphism of K onto a subfield of K.
\end{enumerate}

\begin{enumerate}
\def\labelenumi{(\alph{enumi})}
\item
  Show that \(u(\mathfrak{p}) \subset \mathbf{A}\) and \(u(\mathbf{U}) \subset \mathbf{U}\) . (To prove the first point, observe that for all \(z \in \mathfrak{p}\) the equation \(x^n = 1 + z\) admits a solution in \(\mathbf{A}\) for all \(n > 0\) prime to the characteristic of the residue field of \(v\) , using Hensel's Lemma; if \(v(u(z)) < 0\) , deduce that the integer \(v(u(z))\) would be divisible by every integer \(n > 0\) prime to the characteristic of the residue field of \(v\) . Then deduce the second assertion from the first.)
\item
  Deduce from (a) that, either \(u(\mathbf{K}^*) \subset \mathbf{U}\) , or \(u\) is continuous (consider the image under \(u\) of a uniformizer of \(v\) ).
\item
  Give an example of an algebraically closed field \(\Omega\) such that, taking \(\mathbf{K} = \Omega((\mathbf{X}))\) , \(\mathbf{K}\) is isomorphic to a subfield of \(\mathbf{K}\) contained in \(\mathbf{U} \cup \{0\}\) (cf.~Algebra, Chapter V, § 5, Exercise 13).
\end{enumerate}

¶ 8. (a) Let K be a field, v a valuation on K of height 1 and A its ring. Let H be a compact subset of K (with the topology \(T_{v}\) ) and \(a \neq 0\) a point of K. Show that there exists a polynomial \(f \in K[X]\) without constant term and such that \(f(a) = 1\) , \(f(H) \subset A\) . (Prove that f may be taken to be a polynomial of the form

\[
1 - (1 - a ^ {- 1} \mathrm{X}) (1 - c _ {1} ^ {- 1} \mathrm{X}) ^ {n (1)} \dots (1 - c _ {r} ^ {- 1} \mathrm{X}) ^ {n (r)}
\]

where the \(c_{i}\) are suitably chosen elements of \(\mathbf{H}\) such that \(v(c_{i}) < v(a)\) and the \(n(i)\) are sufficiently large integers \(>0\) .

\begin{enumerate}
\def\labelenumi{(\alph{enumi})}
\setcounter{enumi}{1}
\tightlist
\item
  Let X be a totally disconnected compact space; let the ring \(\mathcal{C}(\mathbf{X};\mathbf{K})\) of continuous mappings from X to K be given the uniform convergence topology. Let B be a subring of \(\mathcal{C}(\mathbf{X};\mathbf{K})\) containing the constants and separating the points of X; show that B is dense in \(\mathcal{C}(\mathbf{X};\mathbf{K})\) (Use (a) and General Topology, Chapter X, §4, Exercise 21 (b).)
\end{enumerate}

¶ 9. Let A be a discrete valuation ring, π a uniformizer of A and K the field of fractions of A; suppose that A is complete with the π-adic topology. The injective A-modules (Algebra, Chapter II, §2, Exercise 11) are identical with the divisible A-modules (Algebra, Chapter VII, §2, Exercise 3) and are direct sums of A-modules isomorphic either to K or to K/A (Algebra, Chapter VII, §2, Exercise 3). Moreover, every monogenous A-module is isomorphic to a submodule of K/A. The A-module Hom \(_{A}\) (M, K/A) is called the (algebraic) toric dual of an A-module M and denoted by M*; the canonical mapping c \(_{M}\) : M → M** is injective (Algebra, Chapter II, §2, Exercise 13 (b)). For every submodule N of M the submodule N \(^{0}\) of M* consisting of the u such that u(x) = 0 is called the orthogonal of N in M*; the dual of M/N is canonically identified with N \(^{0}\) and the dual of N with M*/N \(^{0}\) . The toric dual of a direct limit lim M \(_{\alpha}\) is canonically isomorphic to the inverse limit lim M \(_{\alpha}^{*}\) .

\begin{enumerate}
\def\labelenumi{(\alph{enumi})}
\tightlist
\item
  We know that the A-modules of rank one (Algebra, Chapter VII, § 4,
\end{enumerate}

Exercise 22) are isomorphic to a module of one of the form A, K, K/A or A/ \(\pi^{h}\) A. Show that the algebraic toric duals of K and A/ \(\pi^{h}\) A are respectively isomorphic to K and A/ \(\pi^{h}\) A, that the toric dual of A is isomorphic to K/A and that that of K/A is isomorphic to A (use the knowledge about the sub-A-modules of K and the fact that A is complete (with respect to the dual of K/A)). Deduce that, for every A-module M of finite rank, M* is a module of the same rank and that the canonical homomorphism \(c_{M}\) is bijective.

\begin{enumerate}
\def\labelenumi{(\alph{enumi})}
\setcounter{enumi}{1}
\item
  Let M be an A-module and N a submodule of M of finite rank. Show that the orthogonal of \(N^{0}\) in \(M^{**}\) is identified (by \(c_{M}\) ) with N (use (a)).
\item
  Show that an A-module M which is Noetherian (resp. Artinian) is of finite rank (embed M in its injective envelope (Algebra, Chapter II, §2, Exercise 18)); then M* is Artinian (resp. Noetherian).
\item
  Let \(\sigma(M^{*}, M)\) denote the topology (on \(M^{*}\) ) of pointwise convergence on \(M\) ( \(K/A\) being given the discrete topology). Show that, if \(N\) is a submodule of \(M\) , \(\sigma(N^{0}, M/N)\) is induced on \(N^{0}\) by \(\sigma(M^{*}, M)\) and that \(\sigma(M^{*}/N^{0}, N)\) is the quotient by \(N^{0}\) of the topology \(\sigma(M^{*}, M)\) . (For the second point, note that, if \(P\) is a submodule of \(M\) such that \(N \subset P\) , the dual of \(P/N\) is identified with \(N^{0}/P^{0}\) ). If \(M = \lim M_{\alpha}\) , the topology \(\sigma(M^{*}, M)\) is the inverse limit of the topologies \(\sigma(M_{\alpha}^{*}, \overrightarrow{M_{\alpha}})\) .
\item
  The topologies \(\sigma(\mathbf{K},\mathbf{K})\) and \(\sigma(\mathbf{A},\mathbf{K}/\mathbf{A})\) are the \(\pi\) -adic topologies; the topologies \(\sigma(\mathbf{A}/\pi^{h}\mathbf{A},\mathbf{A}/\pi^{h}\mathbf{A})\) and \(\sigma(\mathbf{K}/\mathbf{A},\mathbf{A})\) are the discrete topologies. Deduce that for every A-module M, the module \(\mathbf{M}^{*}\) with \(\sigma(\mathbf{M}^{*},\mathbf{M})\) is linearly compact (Chapter III, §2, Exercise 15; consider M as a quotient module of a free A-module).
\item
  Let M, N be two A-modules; for every homomorphism \(u: \mathbf{M} \to \mathbf{N}\) , let \(^{t}u\) denote the homomorphism \(\operatorname{Hom}(u, \mathbf{l}_{\mathbf{K/A}})\) from \(\mathbf{N}^{*}\) to \(\mathbf{M}^{*}\) such that
\end{enumerate}

\[
\left(^ {t} u (w)\right) (x) = w (u (x))
\]

for all \(x \in \mathbf{M}\) and all \(w \in \mathbf{N}^*\) ; show that \({}^t u\) is continuous for the topologies \(\sigma(\mathbf{N}^*, \mathbf{N})\) and \(\sigma(\mathbf{M}^*, \mathbf{M})\) . If \(u\) is the endomorphism \(x \mapsto \pi x\) of \(\mathbf{M}\) , \({}^t u\) is the endomorphism \(w \mapsto \pi w\) of \(\mathbf{M}^*\) . For every submodule \(\mathbf{P}\) of \(\mathbf{M}\) , \((u(\mathbf{P}))^0 = {}^t u^{-1}(\mathbf{P}^0)\) .

¶ 10. The hypotheses and notation are the same as in Exercise 9.

\begin{enumerate}
\def\labelenumi{(\alph{enumi})}
\item
  If M is a discrete topological A-module, show that M is a torsion module. If further M is linearly compact, show that M is Artinian (if N is the kernel of the endomorphism \(x \mapsto \pi x\) , observe that N is linearly compact and discrete and can be considered as a vector (A/ \(\pi\) A)-space; then use Exercise 20 (d) of Chapter II, § 2).
\item
  Deduce from (a) that every linearly compact topological A-module is strictly linearly compact (Chapter II, § 2, Exercise 19).
\item
  The topological toric dual of a topological A-module M is the submodule \(M'\) of the algebraic toric dual \(M^{*}\) consisting of the continuous homomorphisms from M to K/A (the latter being given the discrete topology). If the topology on M is linear (Chapter II, § 2, Exercise 14) and N is a closed submodule of M, show that the topological toric dual of M/N is identified with \(N^{0} \cap M'\) and that of N with \(M' / (N^{0} \cap M')\) (to determine the dual of N, note that, for every continuous homomorphism u from N to K/A, there exists an open submodule U of M such that \(u(x) = 0\) in \(N \cap U\) , then use Exercise 9).
\item
  For a topological A-module M of finite rank, the topological toric dual \(M'\) is equal to the algebraic toric dual \(M^{*}\) (reduce it the case of modules of rank one).
\item
  Let M be a Hausdorff topological A-module whose topology is linear; show that for all \(x \neq 0\) in M, there exists \(u \in M'\) such that \(u(x) \neq 0\) (note that there exists an open submodule U of M such that \(x \notin U\) ); in other words, the canonical mapping \(\mathbf{M} \to (\mathbf{M}')^*\) is injective. Deduce that, if N is a closed submodule of M, \(N^0 \cap M'\) is dense in \(N^0\) with the topology \(\sigma(\mathbf{M}^*, \mathbf{M})\) (use (d)), and consequently \(N = M \cap (N^0 \cap M')^0\) in \(M^{**}\) .
\end{enumerate}

Show that M is dense in \((\mathbf{M}^{\prime})^{*}\) with the topology \(\sigma((\mathbf{M}^{\prime})^{*}, \mathbf{M}^{\prime})\) .

\begin{enumerate}
\def\labelenumi{(\alph{enumi})}
\setcounter{enumi}{5}
\item
  Let M be a linearly compact topological A-module; show that the canonical injection \(M \rightarrow (M')^{*}\) is a bijection and that the topology on M (identified with \((M')^{*}\) ) is equal to \(\sigma(M, M')\) . (Observe that, if U is an open submodule of M, M/U is Artinian by (a) and therefore \(U^{0}\) is Noetherian (Exercise 9 (c)) and deduce the identities of the topologies considered on M; complete with the aid of (e).)
\item
  Let M, N be two topological A-modules whose topologies are linear; for every continuous homomorphism \(u \colon \mathbf{M} \to \mathbf{N}\) , \(^t u(\mathbf{N}') \subset \mathbf{M}'\) and \(^t u\) is also used (by an abuse of notation) to denote the linear mapping from \(\mathbf{N}'\) to \(\mathbf{M}'\) with the same graph as \(^t u\) . If M and N are Hausdorff, show that the restriction to M of \(^t (^t u)\) coincides with \(u\) (M and N being respectively considered as canonically embedded in \((\mathbf{M}')^*\) and \((\mathbf{N}')\) ); moreover, for every submodule Q of N,
\end{enumerate}

\[
(\bar {u} ^ {1} (\mathbf {Q})) ^ {0} \cap \mathbf {M} ^ {\prime} = ^ {t} u (\mathbf {Q} ^ {0})
\]

(use (e)).

\begin{enumerate}
\def\labelenumi{(\alph{enumi})}
\setcounter{enumi}{7}
\tightlist
\item
  Let M be a linearly compact A-module; deduce from (g) and Exercise 9 (f) that for M to be torsion-free, it is necessary and sufficient that M' be divisible and that for M to be divisible, it is necessary and sufficient that M' be torsion-free.
\end{enumerate}

